%% Paper 23 -- manuscript, generic article class (production form:
%% anonymous + TMLR style; see DECISION_PACKET.md). Target venue: TMLR.
%% Author: Ranadhir Ghosh (sole author). Drafted 2026-11-10 from the
%% frozen census artefacts in ../results/census/ (Arm A / Arm B JSON,
%% design + audit tables) and ../REPORT.md. Every number in this file
%% traces to those artefacts; the per-number map is in
%% DATA_PROVENANCE.md in this folder.
\documentclass[11pt]{article}

\usepackage[margin=1in]{geometry}
\usepackage{amsmath,amssymb}
\usepackage{booktabs}
\usepackage{graphicx}
\usepackage[hidelinks]{hyperref}
\usepackage{url}
\usepackage{xcolor}

\newcommand{\sT}{\sigma^2_{\mathrm{task}}}
\newcommand{\sR}{\sigma^2_{\mathrm{run}}}
\newcommand{\sG}{\sigma^2_{\mathrm{grader}}}

\title{What a Benchmark Score Cannot Say:\\
Variance Components and Detectability Limits on Three Agent Benchmarks}

\author{Ranadhir Ghosh\\
FIS Global, Jacksonville, FL, USA}

\date{}

\begin{document}
\maketitle

\begin{abstract}
When two agents differ by a few points on a public benchmark, is that a
capability difference or measurement noise? We answer with a
pre-registered measurement study of the three public agent benchmarks
whose repeated runs survive artefact-level ingest verification (94,555
run outcomes; three of six candidates failed verification and were
excluded): we decompose score variance into task-sampling and run-level
components with two co-primary estimators --- method of moments and a
Bayesian hierarchical model --- under a frozen rule for reporting their
disagreement, and we convert the components into per-benchmark
minimum-detectable-effect (MDE) design tables and an audit of published
gaps on the two benchmarks that publish auditable ones. Run-level
variance is 0.38--0.71 of task-sampling variance on the two benchmarks
where the pre-registered verdict is readable, and is indeterminate,
not small, on the third (a two-agent benchmark); the two estimators
disagree on the level by 0.09--0.11 throughout, a disagreement we
report rather than average away, and pooled ratios summarise wide
per-agent spreads. Detectability, not repetition, is the binding
constraint: on the METR task suite a five-point gap cannot be resolved
by any number of runs at the headline design and needs 257 tasks at
eight runs, while both of $\tau$-bench's official leaderboard gaps fall
below the minimum detectable effect of the design that produced them.
On the benchmarks the public record lets anyone measure, a score as
currently reported does not carry the information needed to say
whether the differences it advertises are real.
\end{abstract}

\noindent\textbf{Keywords:} agent evaluation; benchmark metrology;
variance components; minimum detectable effect; reproducibility

\section{Introduction}

Leaderboards rank agents by point scores. A two- or three-point gap
between adjacent entries is routinely read --- in papers, in model
cards, and in procurement decisions --- as a capability difference.
That reading assumes the score is a precise measurement of the agent.
It is not obviously so: an agent benchmark score is a mean over a
finite set of tasks, each attempted a small number of times by a
stochastic system, and every one of those finitudes contributes
variance that the published number does not report. One live industry
leaderboard's methodology page states the position plainly: its
intervals ``reflect within-run case-level sampling, not pass-to-pass
stability'', and ``run-to-run variance is not yet measured for this
board'' \cite{algolia2026}. The field's own systematic survey reaches
the same conclusion at scale: across 259 studies, non-determinism is
named an uncontrolled and unmeasured threat to benchmark conclusions
\cite{nageshwaran2026}.

The components of that unmeasured variance have each been touched
somewhere. Miller \cite{miller2024} gives the reference treatment of
task-sampling uncertainty for static question-answering evaluations,
with paired comparisons and power analysis, but no run, grader, or
environment component and nothing agentic. Bjarnason et al.\
\cite{bjarnason2026} measure run-to-run spread on one benchmark
(single-run pass@1 varies by 2.2--6.0 points across runs) and recommend
power analysis, but decompose nothing and audit nothing. Zhou et al.\
\cite{zhou2026} decompose variance on one static code benchmark into
two components and derive a repetition ceiling; Messing
\cite{messing2026} decomposes the uncertainty of LLM annotation
pipelines, including interactions, but not of repeated agent runs.
The machinery we use is standard experimental-design machinery ---
ANOVA-style variance components in the line of Bouthillier et al.\
\cite{bouthillier2021} and generalizability theory, and power analysis
in the line of Card et al.\ \cite{card2020} and Miller
\cite{miller2024} --- and we claim no novelty for it. What none of the
works above provides --- and what a reader of a leaderboard actually
needs --- is the three-part conjunction: a decomposition of an agent
benchmark's score variance on genuine repeated runs, the conversion of
those components into the smallest gap the benchmark's own design can
detect, and an audit of the published record against that detectability
limit. This paper provides that conjunction, under a pre-registration
frozen before any census data was touched, on every public benchmark
whose repeated runs survive verification at the artefact itself.

The contributions are:

\begin{enumerate}
\item \textbf{A variance census with two co-primary estimators and a
frozen disagreement rule.} Across the three benchmarks that pass
ingest verification, the ratio of run-level to task-sampling variance
is $\rho = 0.3767$ [0.3154, 0.4426] (METR task suite) and $0.7134$
[0.6355, 0.7844] (SWE-bench Verified repeated-run corpus) under the
moment estimator, with $\tau$-bench indeterminate at $0.9834$
[0.7798, 1.2239]. The Bayesian arm reads 0.09--0.11 higher on all
three and is formally discordant with the moment arm throughout; both
arms are reported in full, and the verdicts are read from the arm the
frozen threshold was defined on. Estimator disagreement at this size
is itself a measurement result: it bounds how precisely the field's
current tools can state the noise level of a benchmark.
\item \textbf{Detectability tables, including the regime where runs
cannot help.} Per benchmark, the minimum detectable effect at 80\%
power as a function of runs and tasks, with required-runs and
required-tasks inverses. The headline case is an impossibility, not a
shortfall: where the task term binds, no number of runs resolves the
gap, and the design table must say so.
\item \textbf{An audit of published gaps against their own designs.}
Both gaps on $\tau$-bench's official leaderboard are smaller than the
minimum detectable effect at the design actually used; of 24 published
contrasts on the SWE-bench Verified corpus, 18 clear their design's
MDE, 6 fall below it, and 3 change classification with the estimator.
One benchmark contributes no auditable gap at all, because it
publishes none with a stated design --- a finding about reporting
practice, reported as such.
\end{enumerate}

A note on what this paper is not. It is a measurement of published
benchmark designs and records, not a proposal for a new benchmark, not
a validity or shortcut analysis (a different object: bias, not
variance \cite{zhu2025,wang2026}), and not a re-analysis of any single
benchmark's capability claims. Its negative results --- an
indeterminate benchmark, an empty audit leg, three excluded corpora
--- are reported at the same prominence as its positive ones, because
under a pre-registered design they are part of the answer.

\section{The measurement problem}

A published score for agent $i$ is a mean over a finite design,
$\hat{S}_i = \frac{1}{T}\sum_j \frac{1}{R_j}\sum_r Y_{ijr}$, where $j$
indexes tasks drawn from the benchmark's task set (size $T$), $r$
indexes independent runs (count $R_j$), and $Y_{ijr}$ is the run
outcome --- binary as the primary endpoint, continuous partial credit
as secondary. Under the random-effects model $Y_{ijr} = p_{ij} +
\varepsilon_{ijr}$, with $p_{ij}$ the agent's per-task success
propensity, the law of total variance splits a single run's outcome
variance, for fixed agent $i$, into named components
(Table~\ref{tab:components}). At design level, the variance of the
published number is approximately
$\mathrm{Var}(\hat{S}_i) \approx \sT/T + \sR/(T\cdot R)$ (plus a grader
analogue where identified): tasks and runs buy precision at different
rates, and only the first term responds to a larger benchmark while
only the second responds to repetition. A leaderboard that reports
neither component cannot say which lever its own noise responds to
--- or whether either lever suffices for the gaps it ranks.

\begin{table}[t]
\centering\small
\begin{tabular}{@{}lll@{}}
\toprule
Component & Definition (per agent $i$) & Generated by \\
\midrule
Task sampling $\sT(i)$ & $\mathrm{Var}_j(p_{ij})$ & Which tasks are in the benchmark \\
Run-level $\sR(i)$ & $\mathrm{E}_j[\mathrm{Var}_r(Y\mid i,j)]$ & Agent stochasticity; environment \\
 & & non-determinism where inseparable \\
Grader $\sG(i)$ & $\mathrm{E}[\mathrm{Var}(Y\mid \text{frozen trajectory, re-grading})]$ & Stochastic or model-based grading \\
Environment $\sigma^2_{\mathrm{env}}(i)$ & $\mathrm{E}[\mathrm{Var}(Y\mid \text{frozen policy and grading})]$ & Live-environment drift \\
\bottomrule
\end{tabular}
\caption{The components of a benchmark score's variance. The
environment component is counted within the run-level component
wherever a corpus cannot identify it separately (refinement R1,
Section~\ref{sec:design}); the grader component is estimated only
where a corpus re-grades frozen trajectories, and is otherwise an
estimability limit, never an imputed zero.}
\label{tab:components}
\end{table}

The falsifier statistic is the ratio $\rho = (\sR + \sG)/\sT$, pooled
across agents with equal weights, per benchmark. In every corpus this
census could measure, $\sG$ is zero by construction (programmatic
grading, no re-grades), so the estimated numerator is the run-level
component alone; $\rho$ is written in its general form because the
estimand, frozen before the census, is defined for judged benchmarks
too --- and on judged benchmarks it remains unmeasured, here and, to
our knowledge, everywhere (Section~\ref{sec:limits}). The thesis of
this paper --- that single-run scores are not as reliable as the
field's usage assumes --- dies, by the pre-registered criterion, if
$\rho < 0.10$ on a set of at least three benchmarks with genuine
repeated runs. The threshold is deliberately ungenerous to the thesis:
run noise one-tenth the size of task noise would already vindicate
current practice.

\section{Design: a pre-registered census}
\label{sec:design}

\subsection{Freezing, and what was frozen}

The full design --- estimands, selection rule, falsifier, estimators,
diagnostics, and the disagreement rule below --- was frozen on
2026-10-31 (addendum timestamp 2026-10-31T16:00:07Z) in a SHA-256 hash
chain, before any census data was downloaded, ingested, or computed
on, and before the model-based arm's code existed. The bounded pilot
that preceded it (one agent on one benchmark, $\rho = 0.245$, 95\% CI
[0.134, 0.397]) was itself run under a pre-registration frozen
2026-10-30; its reading is indicative only, since the falsifier's
verdict requires the full set. Two clarifications were approved with
the investigation packet and carried into the freeze verbatim:
\textbf{R1} --- where a corpus cannot separately identify the
environment component, it is counted within the run-level component
and the inseparability reported as an estimability limit; the same
treatment applies to the grader component where no re-grading exists
(it is then either zero by construction, for deterministic
programmatic grading, stated as such, or not identified). Amendments
to any frozen quantity would have required a new frozen amendment
before the affected outcomes were examined; none was made.

\subsection{The census set: selection and ingest verification}

A benchmark enters the set if and only if (a) genuine repeated runs
($n_{ij} \geq 2$ for a substantial share of agent$\times$task cells)
are publicly available at run level, (b) outcomes are binary or
binarisable without judgement calls, and (c) task identities are
stable across runs. Six candidates were carried into the census. Each
was then subjected to \emph{ingest verification}: the outcome-bearing
files must be enumerated and parsed at the artefact itself, their
counts reconciled against the corpus's stated design, and the licence
permitting research reuse confirmed. Repository existence is not data
presence --- a rule this programme learned before the census, when a
headline 2026 result in this exact field advertised a public outcome
tensor that its repository does not contain at any ref (the kill is
recorded in the replication record; that work's numbers are cited in
this paper only as its authors' claims, never as verified).

Three candidates pass (Table~\ref{tab:census}); three fail, and the
failures are findings:

\begin{itemize}
\item \textbf{B2 (HAL generalist traces): FAIL.} The dataset returns
HTTP 401 at census time across its page, file tree, and file-resolve
endpoints; the outcome files (encrypted archives) can be neither
enumerated nor parsed.
\item \textbf{B6 (HAL CORE-Bench): FAIL.} The same access failure;
its frozen ``re-verified at census'' repeat-structure check is
therefore impossible.
\item \textbf{B5 (de la Rubia \& Pafka coding-agent runs): FAIL on the
selection legs.} The corpus is present and MIT-licensed (468 run
records), but its outcome is a continuous holdout AUC with no
pre-stated binary bar, and it contains exactly one task, so the task
component has no population to estimate over.
\end{itemize}

The passing set is exactly the pre-registered floor of three. That
thinness is itself a result about the field's data practices: of six
candidate benchmarks chosen \emph{because} they advertise repeated
runs, half cannot support the measurement, two of them because
advertised-open data is not, at the artefact, open.

\begin{table}[t]
\centering\small
\begin{tabular}{@{}lrrrr@{}}
\toprule
Benchmark & Outcomes & Agents & Tasks & Cells $n\geq2$ \\
\midrule
B1: METR task suite (HCAST + RE-Bench) & 39,363 & 52 & 177 & 0.979 \\
B3: $\tau$-bench (retail + airline) & 1,980 & 2 & 165 & 1.000 \\
B4: Bjarnason et al.\ corpus (SWE-bench Verified) & 53,212 & 12 & 500 & 0.991 \\
\bottomrule
\end{tabular}
\caption{The census set. B4's agents are agent configurations (two
scaffolds $\times$ three models $\times$ two temperatures, not all
cells present). B1 has 51 estimable agents; one (gpt2) has no repeated
cells and is excluded by the $n_{ij}\geq2$ rule. In B1, 108 identical
duplicate run records were counted once and 308 conflicting re-listings
excluded, both logged in the ingest ledger.}
\label{tab:census}
\end{table}

\subsection{Anchors}

Before any variance estimate from a benchmark is trusted, a
published-style score is recomputed from the raw run-level records two
independent ways, with a record-count cross-check. B1's estimates
reproduce the sealed pilot artefact bit-for-bit on overlapping slices
(maximum absolute difference 0.0); B4's per-configuration means
reproduce the source artefact to $1.4\times10^{-14}$; B3's two
published contrasts are reproduced by the census corpus itself (its
trajectories are the leaderboard runs). One discrepancy is on record:
the pilot document's printed slice table is not reproducible from its
sealed artefact; the census anchors on the artefact, and the printed
table stands as an uncorrected reporting error in the frozen document.
The frozen thresholds are design constants and are unaffected.

\subsection{Two co-primary arms, and the disagreement rule}

\textbf{Arm A (method of moments)} estimates the components without
distributional assumptions: within-cell variance
$s^2_{ij} = \frac{n_{ij}}{n_{ij}-1}\hat{p}_{ij}(1-\hat{p}_{ij})$,
averaged with $n_{ij}$ weights for $\sR$; $\sT$ by the standard ANOVA
moment correction (observed spread of cell means minus its expected
sampling inflation, truncated at zero with the truncation reported).
Uncertainty is by cluster bootstrap over tasks ($B = 2{,}000$), with
BCa intervals reported throughout. One methodological deviation is on
record here rather than repaired silently: the frozen registration
specifies percentile bootstrap intervals; the census reports BCa. Both
interval types are deposited in the artefacts; the fresh-tester audit
of this census verified that the percentile intervals yield identical
verdicts and identical disagreement-rule outcomes, so no finding in
this paper depends on the substitution. \textbf{Arm B (Bayesian partial pooling)} fits, per
agent, a logistic random-effects model
$\mathrm{logit}(p_{ij}) = \mu_i + v_{ij}$, $v_{ij} \sim
\mathrm{Normal}(0, \tau_i)$, with priors $\mu_i \sim \mathrm{Normal}(0,
2.5^2)$ and $\tau_i \sim \mathrm{HalfNormal}(0,1)$, and derives the
same probability-scale estimands per posterior draw, so both arms
estimate the same quantities. A continuous-outcome arm (LMM on the
original scale) is secondary in both. Sampler: NUTS, 4 chains, 2,000
warm-up + 2,000 draws, target acceptance 0.95. Diagnostics
($\hat{R} \leq 1.01$, bulk/tail ESS $\geq 400$, zero divergences,
posterior predictive checks) carry frozen consequences: exactly one
remediation at stricter settings, then FAILED, recorded and never
dropped. A pre-specified prior-sensitivity refit (widened priors) is
always reported alongside.

Neither arm may be picked after results are seen, so the freeze
specifies the \textbf{disagreement rule} in advance. The arms
materially disagree on a benchmark if (a) the pooled $\rho$ points
differ by more than 0.05, (b) the two 95\% intervals lie on opposite
sides of the 0.10 falsifier threshold, or (c) a pooled E1/E2 point
estimate differs by a factor outside $[2/3, 1.5]$. On material
disagreement: both arms are reported side-by-side in full, the
benchmark is flagged DISCORDANT, and the formal verdict is read from
Arm A alone --- because the falsifier, its threshold, and R1 were
frozen in moment-estimator terms, and re-basing a frozen threshold
onto a model-based scale would move the bar. Criterion (c), we record
here, turned out not to be evaluable for this census: the frozen Arm B
pooled artefact carries $\rho$ only, with no pooled components, so no
factor is computed rather than estimated post hoc. That is a
limitation of the frozen artefact set, stated rather than repaired.

\subsection{Detectability and the audit}

For two agents compared on the same task set (the paired design
leaderboards actually use), $\mathrm{SE}_{\mathrm{diff}}(R) =
\sqrt{\sigma^2_{\mathrm{task,diff}}/T +
(\sR(A)+\sR(B))/(T\cdot R)}$ and the minimum detectable effect at
two-sided $\alpha = 0.05$, power 0.80 is $\mathrm{MDE}(R) = 2.802 \cdot
\mathrm{SE}_{\mathrm{diff}}(R)$. The inverse, $R_{\mathrm{required}}
(\delta)$, has an honest edge case: if its denominator is
non-positive, \emph{no number of runs suffices} and the table reports
$T_{\mathrm{required}}$ at the benchmark's actual $R$ instead. That
edge case is a result, not an error. Arm B's MDE construction is
unpaired (its frozen model carries no per-pair estimand); both arms'
tables are reported, and rows whose audit classification flips between
arms are flagged arm-sensitive, never resolved by choosing an arm.
The audit (E6) classifies each published pairwise gap that can be
sourced to the benchmark's own publication or official leaderboard ---
with the gap, and the $T$ and $R$ of the design actually used,
recorded --- against the MDE at that design. No gap is reconstructed,
estimated, or invented where no published value can be sourced. The
audit is deep rather than broad by design: it asks of each gap the
strongest question (does it clear its own benchmark's detectability
limit?) over few benchmarks, rather than a weaker question over many.

\section{Results}

\subsection{The falsifier verdicts}

Table~\ref{tab:verdicts} gives both arms, side-by-side as the frozen
rule requires. The formal verdicts, read from Arm A: \textbf{B1
SUPPORTED} ($\rho = 0.3767$ [0.3154, 0.4426]); \textbf{B3
INDETERMINATE} ($\rho = 0.9834$ [0.7798, 1.2239]); \textbf{B4
SUPPORTED} ($\rho = 0.7134$ [0.6355, 0.7844]). No benchmark is
refuted: wherever the verdict is readable, run-level variance is
between roughly two-fifths and seven-tenths of task-sampling variance
--- four to seven times the level at which the thesis would have died.
B1's continuous secondary endpoint corroborates its binary primary
($\rho = 0.4350$ [0.3657, 0.5102]).

\begin{table}[t]
\centering\small
\begin{tabular}{@{}lllll@{}}
\toprule
Benchmark & Arm A $\rho$ (BCa 95\%) & Arm A & Arm B $\rho$ (95\% CrI) & Arm B \\
\midrule
B1 (METR) & 0.3767 [0.3154, 0.4426] & SUPPORTED & 0.4709 [0.4537, 0.4886] & SUPPORTED$^\dagger$ \\
B3 ($\tau$-bench) & 0.9834 [0.7798, 1.2239] & INDET. & 1.0823 [0.9200, 1.2741] & INDET. \\
B4 (Bjarnason) & 0.7134 [0.6355, 0.7844] & SUPPORTED & 0.8237 [0.7960, 0.8526] & INDET.$^\ddagger$ \\
\bottomrule
\end{tabular}
\caption{Falsifier verdicts, both arms. $^\dagger$B1 Arm B carries the
prior-sensitivity flag (Section~\ref{sec:discord}). $^\ddagger$B4 Arm
B's indeterminacy is a fit-failure verdict under the frozen rule
(Section~\ref{sec:fits}), not an evidence verdict; its pooled point
covers the 9 agents whose primary fits passed. Formal verdicts are
read from Arm A alone under the frozen disagreement rule; all three
benchmarks are DISCORDANT on criterion (a).}
\label{tab:verdicts}
\end{table}

B3's indeterminacy deserves plain statement, because it is the kind of
result a less disciplined write-up would round into support. With two
agents, the pooled estimate rests on a cross-fitted pair ratio; its
point estimates sit at or above $\rho = 1$ --- run variance comparable
to or larger than task-sampling variance --- but the frozen SUPPORTED
clause requires the interval to lie entirely inside the definitional
range (0, 1), which intervals crossing 1 fail. The frozen rules
therefore record indeterminate. That is a precision statement, not a
weak effect, and it is also a statement about $\tau$-bench's published
design: a two-agent repeated-run corpus cannot certify the quantity
its leaderboard most needs. Per-agent readings underline the spread
the pooled numbers summarise: within B1, per-agent $\rho$ runs from
0.15 to 15.6 (median 0.35; the maximum is a near-degenerate
low-task-variance agent, reported in the artefacts, not trimmed), and
B3's two agents read 0.93 and 1.05 individually.

The grader component is an estimability limit throughout: no corpus
in the set re-grades frozen trajectories, and grading in all three is
programmatic, so E3 is zero by construction in all three --- stated,
not estimated. Nothing in this census measures grader variance on
judged benchmarks; Section~\ref{sec:limits} returns to this.

\subsection{The disagreement, reported}
\label{sec:discord}

The disagreement rule fired on criterion (a) for all three benchmarks:
$|\Delta\rho| = 0.0942$ (B1), $0.0989$ (B3), $0.1103$ (B4). Criterion
(b) fired nowhere --- every interval in both arms lies entirely above
the 0.10 threshold, so the arms never disagree about which side of the
falsifier the evidence sits on. They disagree about the level, by
about a tenth of the task component, consistently in one direction
(the model arm reads higher). We do not resolve this. The
registration names candidate causes --- sparse cells, prior influence
at the sensitivity refit, model misfit --- and forbids post-hoc
resolution; what can be said from the record is where each flag sits.
B1's Arm B carries the prior-sensitivity flag: widening the priors
moves its pooled $\rho_B$ by 0.0748, above the 0.05 flag threshold.
The caveat that travels with that flag is that the widened-prior pool
covers only 28 of 51 agents (23 sensitivity fits failed the frozen
diagnostics), so the shift mixes a prior effect with an agent-set
effect, and the frozen construction does not decompose the two. B3's
corresponding shift is 0.0499 --- below the flag by 0.0001, recorded
as borderline, not flagged. B4's is 0.0015, no flag. For a reader who
wants one number: the honest one-number summary of this census is not
a number but a range with a rule --- the moment arm's reading, the
model arm's reading 0.09--0.11 higher, and the pre-registered reason
the formal verdict follows the former.

\subsection{Fit accounting}
\label{sec:fits}

Arm B's frozen list is 297 fits. Final tally: \textbf{246 PASS, 12
REMEDIATED\_PASS, 39 FAILED, 0 FLAGGED}. Every failure followed the
frozen ladder exactly and is recorded in its fit file with
first-attempt and remediation diagnostics. Two clusters matter. In
B1, 23 of 51 widened-prior sensitivity fits failed (Section~\ref{sec:discord}),
and 2 of 51 primary fits failed (davinci\_002,
gpt\_3\_5\_turbo\_instruct --- neither a headline agent; 3.9\% $<$
25\%, so the Arm B verdict stands; misfit share 0.0196). In B4, the
primary fits for \emph{both} headline-pair configurations
(Devstral-2 at temperatures 0 and 0.2) and for Qwen3-32B at
temperature 0 failed on bulk effective sample size (359.6 falling to
277.7 after remediation, against the 400 bar; $\hat{R} \approx$
1.005--1.007; zero divergences throughout). Under the frozen
headline-failure rule, B4's Arm B verdict is therefore indeterminate
regardless of its pooled point. We regard the refusal to rescue these
fits --- no third attempt, no relaxed diagnostic, no model change ---
as part of the method, not an execution shortcoming: a metrology
paper that quietly re-tunes its own instrument until the readings
agree would be making the field's mistake at one remove.

\subsection{Detectability: what the designs can resolve}

Table~\ref{tab:mde} gives the moment arm's paired-design MDEs for each
benchmark's headline pair; the model arm's unpaired tables (uniformly
larger, since they cannot exploit shared tasks) are in the deposited
artefacts and summarised in the text where they change a conclusion.

\begin{table}[t]
\centering\small
\begin{tabular}{@{}lccccccl@{}}
\toprule
Benchmark (headline $T$) & $R{=}1$ & $R{=}2$ & $R{=}4$ & $R{=}8$ & $R_{\mathrm{actual}}$ & $\delta{=}0.05$ needs & $\delta{=}0.02$ needs \\
\midrule
B1: METR ($T{=}163$) & 0.1045 & 0.0832 & 0.0702 & 0.0627 & 0.0605 ($R{\approx}10.9$) & 257 tasks @ $R{=}8$ & 1,603 tasks @ $R{=}8$ \\
B3: $\tau$-bench ($T{=}165$) & 0.1243 & 0.0981 & 0.0820 & 0.0725 & 0.0758 ($R{=}6$) & 348 tasks @ $R{=}8$ & 2,171 tasks @ $R{=}8$ \\
B4: Bjarnason ($T{=}500$) & 0.0461 & 0.0328 & 0.0236 & 0.0171 & 0.0156 ($R{\approx}9.9$) & $R{=}1$ & $R{=}6$ \\
\bottomrule
\end{tabular}
\caption{Minimum detectable effects (score points, normal-critical
values; $t$-critical values are within 0.001 throughout and change no
verdict) for the headline pair of each benchmark, paired design, Arm
A. ``Needs'' columns give the binding resource: where runs cannot
reach $\delta$ at the headline $T$, the required task count at $R=8$;
where they can, the required runs at the headline $T$. A 10-point gap
needs only $R = 2$ on B1 and B3.}
\label{tab:mde}
\end{table}

Read B1's row as the paper's central practical claim. Going from one
run to the $\approx$11 runs actually used buys a factor of 1.7 in
detectability (0.1045 to 0.0605) --- and then the curve is spent: at
$R = 8$ the MDE is 0.0627, and no further repetition moves it
materially, because the remaining variance is the task term, which
repetition does not touch. A five-point gap --- larger than most
adjacent leaderboard gaps --- is simply not resolvable on this design;
resolving it is a task-set decision (257 tasks at $R = 8$), not a
compute decision. The statement is a consequence of the measured
components under the stated paired design model, not an empirical
impossibility claim about METR scores in general; and because the two
arms' MDE constructions differ (Arm A's is paired, Arm B's unpaired),
their tables are not two estimates of one quantity --- Arm B's larger
values are what the same components imply when shared tasks cannot be
exploited, and the audit reports both rather than blending them. B3 is worse in the same way (348 tasks needed, on a
165-task benchmark). B4 inverts the pattern for an instructive reason:
its headline pair is one model at two temperatures, whose
task-difference variance is 0.0020 --- two orders below B1's 0.0609 ---
so runs bind and small gaps are reachable ($R = 6$ for two points).
Detectability is a property of the pair and the design, not of the
benchmark's size or budget; it has to be computed, per comparison, from
measured components. At present it is computed by no one.

\subsection{The audit}

\textbf{B3 (2 rows).} $\tau$-bench's official leaderboard reports
pass$^1$ for gpt-4o and claude-3.5-sonnet on retail (0.604 vs 0.692,
gap 0.088) and airline (0.420 vs 0.460, gap 0.040), from designs of
$T = 115$, $R = 4/8$ and $T = 50$, $R = 4/8$ respectively (the
leaderboard values equal the census corpus's own pass$^1$: these are
the leaderboard runs). The MDEs at those designs are 0.0927 and
0.1412. \textbf{Both published gaps are below the minimum detectable
effect of the design that produced them} --- under the moment arm, and
identically under the model arm (MDEs 0.1383 and 0.2097), so the
classification is not arm-sensitive. The retail gap misses its MDE by
0.005: close enough that a slightly larger trial count would have
resolved it, which is precisely the kind of statement a design table
exists to make before, not after, publication.

\textbf{B4 (24 rows).} Of the 24 contrasts in the source paper's
Table 1, at the design actually used ($T = 500$, $R = 10$): 18 are
above their MDE under both arms; 6 are below under both --- four
near-zero temperature contrasts (gaps 0.000--0.016 against MDEs
0.0148--0.0202), a same-temperature model pair on one scaffold (0.005
vs 0.0359), and a cross-scaffold gap (0.012 vs 0.0195); and 3 are
arm-sensitive, above under the moment arm and below under the model
arm (gaps 0.030--0.040). The corpus's headline claims --- model
differences of 0.43--0.47 --- clear their MDEs by an order of
magnitude. A benchmark whose own authors publish repeated runs is, on
this evidence, a benchmark whose headline claims mostly survive their
measurement uncertainty; the failures concentrate exactly where a
reader would hope an audit looks: in the small contrasts a table
invites but its design cannot support.

\textbf{B1 (0 rows).} No pairwise gap could be sourced to METR's own
publications with both a numeric value and the design's $(T, R)$:
METR publishes time horizons and bucketed success rates, not overall
pairwise gaps at a stated design. No gap was reconstructed (the frozen
rule forbids it). The empty leg is reported as a finding about
reporting practice: the benchmark with the richest public run-level
record publishes no gap in a form whose detectability can even be
assessed.

\section{Discussion}

\textbf{What changes for a benchmark maintainer.} Three practices,
each a direct consequence of a table above. (1) Publish the components,
or the run-level records from which they are computed: a score without
its $\sT$ and $\sR$ is a measurement without an instrument reading.
(2) State the MDE next to the gaps a leaderboard advertises, at the
design actually used; the $\tau$-bench retail gap shows the cost of
not doing so is not hypothetical. (3) When a target gap is set ---
``we want to resolve three points'' --- compute which resource binds
before spending on the other: on two of our three benchmarks the
binding resource is tasks, and repetition budgets spent past $R
\approx 4$--$8$ buy almost nothing.

\textbf{What changes for a reader.} A gap below its design's MDE is
not evidence of equality; it is an absence of evidence, and the two
readings license opposite decisions. The audit's arm-sensitive rows
add a second caution: for gaps within a factor of about two of the
MDE, the classification itself can depend on the estimator, which is
an argument for reporting both estimators' design tables --- or for
treating near-MDE gaps as unresolved by construction.

\textbf{Relation to the nearest work.} Zhou et al.\ \cite{zhou2026}
reach a kindred conclusion --- repetition buys little --- on one
static code benchmark, via a two-component generalizability
decomposition and a repetition ceiling. Our design tables say why, in
a form that survives the change of setting: repetition buys little
\emph{when the task term binds}, and whether it binds is a measurable
property of the pair (B4's temperature pair is the counterexample
where repetition is exactly the right spend). We could not verify
Zhou et al.'s numbers: their advertised outcome tensor is absent from
their repository at every ref, and we cite their results as their
claims, with the data-availability finding on record. That finding
is, in miniature, this paper's ingest rule and its B2/B6 exclusions
writ large: in this field, at this moment, \emph{advertised openness
is not evidence of openness}, and any metrology that skips
artefact-level verification inherits the error. Messing
\cite{messing2026} shows interaction terms are real in LLM pipeline
uncertainty; we do not model facet interactions, and record that as a
limitation rather than absorbing it silently. Miller
\cite{miller2024} remains the reference for the task component alone;
the present results are what that framework cannot see, because its
designs have no repeated runs to see it with.

\section{Limitations}
\label{sec:limits}

\begin{enumerate}
\item \textbf{Three benchmarks, at the floor.} The passing set is the
pre-registered minimum. The benchmarks differ in kind (a task suite, a
tool-agent benchmark, a code benchmark's repeated-run corpus), which
supports the conjunction's generality argument, but three is three:
claims are stated per benchmark and the cross-benchmark claim is the
narrow one --- the falsifier's set-level reading --- not a universal
law of agent benchmarks.
\item \textbf{One benchmark has two agents.} B3's pooled estimate is
intrinsically fragile (a cross-fitted pair ratio), and its
indeterminacy is partly a consequence. We report it because excluding
a passing benchmark after seeing its outcomes would violate the
freeze; the same discipline that admits it also caps what it can
certify.
\item \textbf{The grader component is unmeasured.} E3 is zero by
construction on programmatic grading. On judged benchmarks the grader
component could be large (judge biases are documented
\cite{zheng2023}); identifying it requires re-grading frozen
trajectories with judge models --- new inference, separately priced,
explicitly out of this study's registered scope.
\item \textbf{No facet interactions.} Following Messing
\cite{messing2026}, interactions between facets (e.g.\ prompt
$\times$ judge) may contribute variance our additive split attributes
to the named components.
\item \textbf{Arm B's MDE is unpaired.} The frozen model arm has no
per-pair estimand, so its design tables cannot exploit shared tasks
and run uniformly larger; audit rows near the boundary are therefore
arm-sensitive in a way a paired model arm would partly resolve. We
report the asymmetry rather than repair it post hoc.
\item \textbf{Disagreement criterion (c) was not evaluable} (frozen
artefact carries pooled $\rho$ only), and B1's prior-sensitivity shift
confounds prior and agent-set effects. Both are stated where the
numbers are used.
\item \textbf{Protocol deviations, declared.} Three are on the census
record (repaired 2026-11-10; full entries in the census report's
deviation section). (i) The census executed under a newer software
stack (PyMC 5.25.1 and companions) than the frozen lockfile pins (PyMC
5.16.2) after host rebuilds, and the deviation was not logged at the
time; formal verdicts are read from Arm A, which is stack-insensitive
in substance, but exact bit-reproduction of the Arm B fits from the
deposited record is not guaranteed. (ii) Arm A reports BCa intervals
where the registration specifies percentile bootstrap (declared in
Section 3.3; verdicts identical under both). (iii) The SUPPORTED /
INDETERMINATE row labels rest on a reporting taxonomy (the interval
must lie entirely inside $(0, 1)$) that is not itself in the frozen
registration, which freezes only the set-level falsifier; the taxonomy
was applied consistently to all six arm--benchmark rows, and the
set-level falsifier reading --- no benchmark refutes the thesis ---
does not depend on it.
\item \textbf{Static corpora.} HAL's corpus is paused and METR's
versions evolve; the findings are statements about the benchmarks'
published designs and records, framed as such. B4's corpus is a
third-party repeated-run release over SWE-bench Verified tasks, not
the benchmark's own leaderboard runs; its audit rows are that paper's
published contrasts, labelled as such.
\end{enumerate}

\section{Conclusion}

On the public record as it stands, the noise in an agent benchmark
score is measurable, is first-order, and is mostly unreported. Run
noise runs at two-fifths to seven-tenths of task noise where the
verdict is readable; the designs in use cannot resolve the gaps they
are used to advertise, and the binding constraint is usually the task
set, not the run count; and of the benchmarks that advertise the data
needed to check any of this, half fail verification at the artefact.
None of these conclusions required new experiments --- only the
decision to treat a benchmark score as a measurement, and to hold its
reporting to a measurement's standard. The instruments are cheap,
frozen in advance, and deposited with this paper; the design tables
are the deliverable a maintainer can act on Monday morning.

\section*{Data and code availability}

All corpora are public research artefacts: the METR public run-level
records; the $\tau$-bench official release (MIT); the Bjarnason et
al.\ repeated-run release (CC-BY-4.0, Hugging Face
\texttt{ASSERT-KTH/agentic-evals-artifacts}, pinned revision). The
census code, the frozen pre-registration with its hash chain, the
ingest ledger, all 297 fit records, and the JSON artefacts from which
every table in this paper is drawn are deposited in the public
replication package at
\url{https://github.com/ranadhir19/agent-benchmark-metrology-census}
(its root \texttt{validate.py} recomputes the headline census numbers
from the deposited artefacts). The HAL corpora (B2/B6) could not be
redistributed or read; their exclusion is documented in the ingest
ledger.

\section*{Declaration of competing interest}

The author declares no competing interests. This research received no
external funding.

\section*{AI use disclosure}

The author used Muse (Anthropic) for language editing and coding
assistance in the preparation of this manuscript and its deposited
artifacts. All AI-assisted material was reviewed and validated by the
author, who takes full responsibility for the manuscript's content,
accuracy, originality, and integrity.

\begin{thebibliography}{99}

\bibitem{kapoor2024} S.~Kapoor, B.~Stroebl, Z.~S.~Siegel, N.~Nadgir,
A.~Narayanan, ``AI Agents That Matter,'' arXiv:2407.01502, 2024.

\bibitem{nageshwaran2026} V.~Nageshwaran, S.~Ezekiel, T.~T.~Tran,
V.~Lakshmi Narasimhan, ``Large language model agent evaluation and
benchmarking: a systematic survey, meta-taxonomy, and critical
research roadmap,'' \emph{Artificial Intelligence Review}, 2026.
(Bibliographic details as recorded in the programme's verified
literature review; DOI not verified.)

\bibitem{gjorgjevikj2026} A.~Gjorgjevikj, B.~Korousic Seljak,
T.~Eftimov, ``From Scores to Evidence: Robust Evaluation of Foundation
Models,'' tutorial, NeurIPS 2026.

\bibitem{bjarnason2026} B.~H.~Bjarnason, A.~Silva, M.~Monperrus, ``On
Randomness in Agentic Evals,'' arXiv:2602.07150, 2026.

\bibitem{delarubia2026} M.~de la Rubia, B.~Pafka, ``Identical Runs,
Different Results: Benchmarking AI Coding Agents on Open-Weight
Models,'' arXiv:2609.33812, 2026. (As recorded in the programme's
literature review.)

\bibitem{rabanser2026} S.~Rabanser, S.~Kapoor, P.~Kirgis, L.~Liu,
S.~Utpala, A.~Narayanan, ``Towards a Science of AI Agent
Reliability,'' arXiv:2602.16666, 2026.

\bibitem{reliabilitygap2026} ``The Reliability Gap: Repeated-Trial
Consistency and the Deployable Autonomy Horizon of Frontier AI
Agents,'' Research Square preprint rs-11447587/v1, 2026. (Author list
not verified.)

\bibitem{miller2024} E.~Miller, ``Adding Error Bars to Evals: A
Statistical Approach to Language Model Evaluations,''
arXiv:2411.00640, 2024.

\bibitem{bowyer2025} S.~Bowyer, L.~Aitchison, D.~Ivanova, ``Position:
Don't Use the CLT in LLM Evals With Fewer Than a Few Hundred
Datapoints,'' ICML 2025; arXiv:2503.01747.

\bibitem{bouthillier2021} X.~Bouthillier et al., ``Accounting for
Variance in Machine Learning Benchmarks,'' MLSys 2021;
arXiv:2103.03098.

\bibitem{agarwal2021} R.~Agarwal, M.~Schwarzer, P.~S.~Castro,
A.~Courville, M.~G.~Bellemare, ``Deep Reinforcement Learning at the
Edge of the Statistical Precipice,'' NeurIPS 2021.

\bibitem{card2020} D.~Card, P.~Henderson, U.~Khandelwal, R.~Jia,
K.~Mahowald, D.~Jurafsky, ``With Little Power Comes Great
Responsibility,'' EMNLP 2020; arXiv:2010.06595.

\bibitem{kapoor2025hal} S.~Kapoor et al., ``Holistic Agent
Leaderboard: The Missing Infrastructure for AI Agent Evaluation,''
arXiv:2510.11977, 2025.

\bibitem{yao2024} S.~Yao, N.~Shinn, P.~Razavi, K.~Narasimhan,
``$\tau$-bench: A Benchmark for Tool-Agent-User Interaction in
Real-World Domains,'' arXiv:2406.12045, 2024.

\bibitem{kwa2025} T.~Kwa et al., ``Measuring AI Ability to Complete
Long Software Tasks,'' arXiv:2503.14499, 2025.

\bibitem{siegel2024} Z.~S.~Siegel, S.~Kapoor, N.~Nadgir, B.~Stroebl,
A.~Narayanan, ``CORE-Bench: Fostering the Credibility of Published
Research Through a Computational Reproducibility Agent
Benchmark,'' arXiv:2409.11363, 2024.

\bibitem{zhu2025} K.~Zhu et al., ``Establishing Best Practices for
Building Rigorous Agentic Benchmarks,'' arXiv:2507.02825, 2025.
(First author as recorded in the programme's literature review.)

\bibitem{wang2026} B.~Wang, F.~Bianchi, et al., ``Automated Benchmark
Auditing for AI Agents and Large Language Models,''
arXiv:2605.26079, 2026.

\bibitem{zheng2023} L.~Zheng et al., ``Judging LLM-as-a-Judge with
MT-Bench and Chatbot Arena,'' arXiv:2306.05685, 2023.

\bibitem{algolia2026} Algolia, ``LLM Leaderboard --- Methodology,''
\url{https://www.algolia.com/llm-leaderboard/methodology}, accessed
2026-10-30.

\bibitem{zhou2026} Y.~Zhou, J.~Choi, J.~Yao, et al., ``How Much Can
Repeated Running Buy? A Measure Variance Ceiling for LLM Code
Evaluation,'' arXiv:2606.00920, 2026. (Cited as its authors' claims;
the advertised outcome tensor is absent from the public repository ---
see Section 3.2.)

\bibitem{messing2026} S.~Messing, ``Hidden Measurement Error in LLM
Pipelines Distorts Annotation, Evaluation, and Benchmarking,''
arXiv:2604.11581, 2026. (First author as recorded in the programme's
literature review.)

\bibitem{demsar2006} J.~Dem\v{s}ar, ``Statistical Comparisons of
Classifiers over Multiple Data Sets,'' \emph{Journal of Machine
Learning Research} 7 (2006) 229--268.

\bibitem{longjohn2026} J.~Longjohn, V.~Gopalan, et al.,
``Statistical Uncertainty Quantification for Aggregate Performance
Metrics in Machine Learning Benchmarks,'' arXiv:2501.04234. (Venue
details as recorded in the programme's dossier.)

\end{thebibliography}

\end{document}
